\documentclass[11pt,a4paper]{article}

\usepackage[utf8]{inputenc}
\usepackage[T1]{fontenc}
\usepackage{lmodern}
\usepackage[margin=1in]{geometry}
\usepackage{amsmath,amssymb,amsthm,mathtools}
\usepackage{graphicx}
\usepackage{booktabs}
\usepackage{hyperref}
\usepackage{xcolor}
\usepackage{siunitx}
\usepackage{enumitem}
\usepackage{caption}

\hypersetup{
  colorlinks=true,
  linkcolor=blue!50!black,
  citecolor=blue!50!black,
  urlcolor=blue!50!black
}

\title{\bfseries Systroph\=e II: \\
Quantum Field Theory on a Tipler-Pair Background}
\author{Christian Knopp\\\texttt{cknopp@gmail.com}}
\date{2026-05-10}

\begin{document}
\maketitle

\begin{abstract}
We extend the open-source \textsc{Systroph\=e} package
(\texttt{github.com/Zynerji/systrophe}) from a classical-GR
Tipler-pair simulator (v0.1--v0.6) to a quantum-field-theoretic
laboratory (v0.7--v0.13). The new infrastructure supports:
(i) renormalised stress-energy tensor $\langle T_{\mu\nu} \rangle$ on
the Lewis--Papapetrou exterior via point-splitting and the Hadamard
parametrix, with exact recovery of the conformal anomaly
$K_{\text{Kretsch}} / (2880 \pi^2)$ in vacuum; (ii) closed-form
Atiyah--Patodi--Singer $\eta$-invariant on the $\mathbb{Z}_3$
M\"obius cover with verified $\eta_0 + \eta_1 + \eta_2 = 0$
anomaly closure and Callan--Harvey inflow balance; (iii) joint
Floquet quasi-energies on the $(t, b)$ Brillouin product torus with
$\mathbb{Z}_3$ cyclic-permutation symmetry; (iv) Unruh acoustic-
metric mapping ($c^2 - v^2 = F$) identifying the chronology horizon
as an acoustic horizon, with the gravitational and acoustic Hawking
temperatures coinciding to machine precision; (v) Brown--Maclay
$\langle T_{\mu\nu} \rangle$ at a Casimir cavity in flat space with
LP curvature correction; (vi) Newton--Kantorovich back-reaction
solver demonstrating the linear-walk failure of na\"ive Picard
iteration. The package now contains 406 passing tests and ships
the math via a clean module-per-construction layout. We document
the open speculative items as ansatz-level interpretations needing
external input (cavity geometry, vacuum-selection criterion, etc.),
and propose specific experimental tests in BEC sonic-horizon
apparatuses.
\end{abstract}

\section{Introduction}

\textsc{Systroph\=e} v0.1 implemented the closed-form van Stockum /
Tipler exterior with log-periodic structure and a co-axial pair
construction allowing tunable CTC interference. The first whitepaper
\cite{Knopp2026Systrophe} established the classical-GR mathematics:
exact $\alpha = \sqrt{4a^2-1}$, $F(r) \propto \cos(\alpha \ln r/R + \delta)$
sinusoid, phasor-sum closure, and the structural $\delta = \pi$
extinction. The present paper extends to QFT on this background.

\subsection{Scope of Systroph\=e II}

The v0.7--v0.13 series of releases adds the following independent
machinery, each with a clean module and a separate test suite:

\begin{itemize}
\item \textbf{v0.7}: Radial Dirac operator on the LP exterior with
bound-state spectrum
(\texttt{dirac.py}, \texttt{dirac\_spectrum.py}).

\item \textbf{v0.8}: Quantum-field-theory in curved spacetime
infrastructure: Dirac sea pressure, particle creation at the
chronology horizon, QFTCS back-reaction trace, vacuum-energy proxy
(\texttt{dirac\_sea.py}, \texttt{particle\_creation.py},
\texttt{qftcs\_backreaction.py}, \texttt{quantum\_diagnostics.py}).

\item \textbf{v0.9}: Point-splitting renormalisation; full Riemann
tensor, Kretschmann scalar, exact 4D conformal-anomaly trace, and
DeWitt heat-kernel $a_2$ coefficient
(\texttt{point\_splitting.py}).

\item \textbf{v0.11}: Research-grade adiabatic Floquet on the radial
Dirac problem (\texttt{floquet.py}).

\item \textbf{v0.12}: Topological-Casimir mode sums on the
$\mathbb{Z}_3$ cover with the Hurwitz $\zeta$ at $s=-3$ closed
form (\texttt{casimir.py}).

\item \textbf{v0.13}: \emph{This paper.} Eight new modules covering
the post-anomaly machinery (Sections 2--5 below):
\texttt{hadamard\_offtrace}, \texttt{anomaly\_inflow},
\texttt{floquet\_mobius}, \texttt{acoustic\_metric},
\texttt{casimir\_throat}, \texttt{newton\_kantorovich},
\texttt{tipler\_fractal}, \texttt{horned\_torus}; plus four extension
modules \texttt{back\_reaction}, \texttt{floquet\_engineering},
\texttt{dsi\_observables}, \texttt{adm\_export}, \texttt{d\_ctc}.

\item \textbf{v0.15}: Six validation modules promoting previously
ansatz-level claims to concrete calculations (Section 8 below):
\texttt{wormhole\_throat}, \texttt{exotic\_matter\_accounting},
\texttt{chronology\_protection}, \texttt{wilson\_loop},
\texttt{dynamical\_casimir}, \texttt{vacuum\_states}.
\end{itemize}

\section{Renormalised stress tensor and off-trace components}

The package's exact 4D conformal anomaly is
\begin{equation}
\langle T^\mu_{\ \mu} \rangle_{\text{ren}}
  = \frac{K_{\text{Kretsch}}}{2880\, \pi^2}
  \qquad (\text{vacuum},\ R_{\mu\nu} = 0).
\label{eq:anomaly}
\end{equation}
Module \texttt{hadamard\_offtrace.py} extends this to the full
\emph{tensor}
\begin{equation}
\langle T_{\mu\nu} \rangle_{\text{ren}}
  = \frac{1}{2880\, \pi^2}
    R_{\mu\rho\sigma\tau} R_\nu^{\ \rho\sigma\tau},
\label{eq:offtrace}
\end{equation}
whose trace recovers (\ref{eq:anomaly}) exactly. The
state-independent geometric piece (\ref{eq:offtrace}) is testable
against the trace anomaly numerically; state-dependent corrections
(Boulware vs.\ Hartle--Hawking) are documented but not implemented
(they require mode-sum machinery beyond this paper).

We verify the trace-matching condition numerically:
\begin{verbatim}
trace_decomp = trace_decomposition(vs, r=2.0)
trace_decomp["trace_anomaly_check"]  # < 1e-13
\end{verbatim}

\section{Anomaly inflow on the $\mathbb{Z}_3$ cover}

The angular Dirac operator on $S^1$ with twist $\alpha$ has
APS $\eta$-invariant
\begin{equation}
\eta(\alpha) = 1 - 2\alpha,\qquad \alpha \in (0, 1),
\end{equation}
with $\eta(0) = 0$ (symmetric regularisation).

For the $\mathbb{Z}_3$ M\"obius cover ($\alpha_b = b/3 +
\gamma_{\text{eff}}/(2\pi)$), the three branches have
\begin{equation}
(\eta_0, \eta_1, \eta_2)\big|_{\gamma=0} = (0,\ \tfrac{1}{3},\ -\tfrac{1}{3}),
\end{equation}
and
\begin{equation}
\sum_{b=0}^{2} \eta_b = 0,
\label{eq:closure}
\end{equation}
the anomaly closure condition. A non-zero gauge twist
$\gamma_{\text{eff}}$ breaks (\ref{eq:closure}); the residual is
cancelled by a Chern--Simons bulk inflow with coefficient
$k_{\text{CS}} = 1 / (24 \pi^2)$:
\begin{equation}
k_{\text{CS}} \cdot q_{\text{flux}}
  = -\tfrac{1}{2} \sum_b \eta_b(\gamma_{\text{eff}}).
\end{equation}
This is the Callan--Harvey relation evaluated explicitly on the
$\mathbb{Z}_3$ cover; verified by direct construction in
\texttt{z3\_anomaly\_inflow\_balance}.

\section{Acoustic-metric mapping and analog Hawking}

Unruh's 1981 acoustic-metric mapping identifies the LP angular
metric as an acoustic line element via
\begin{equation}
\rho_{\text{acoustic}} = \sqrt{L},\quad
v_\phi = K/\sqrt{L},\quad
c^2 - v_\phi^2 = F.
\end{equation}
The chronology horizon $F = 0$ coincides exactly with an
\emph{acoustic} horizon ($c = v$); the CTC region ($F < 0$) is the
\emph{supersonic} flow region. The acoustic Hawking temperature
\begin{equation}
T_H^{\text{ac}} = \kappa_{\text{ac}} / (2\pi),\qquad
\kappa_{\text{ac}} = \tfrac{1}{2} |dF/dr|_{r=r_H}
\end{equation}
equals the gravitational Hawking temperature at the LP chronology
horizon to machine precision (relative difference $< 10^{-12}$;
\texttt{compare\_acoustic\_vs\_gravitational\_T\_H}).

\paragraph{Experimental access.} The acoustic-analog mapping is the
\emph{only currently lab-tractable path} to the Systroph\=e physics.
In a Bose--Einstein condensate (BEC) with a quantum vortex of
winding $\ell$, the azimuthal flow $v_\phi(r) = \hbar / (m r)$
automatically violates $v_\phi(R) > c/2$ once $R < 2 \xi_{\text{BEC}}$
(the healing length), reproducing the Tipler supercritical
condition $\omega R > 1/2$. A $\mathbb{Z}_3$ triple-vortex
configuration realises the M\"obius cover. The full design
proposal is in \texttt{docs/EXPERIMENTAL\_ACOUSTIC\_ANALOG.md}.

\section{Floquet engineering and discrete-scale invariance}

\subsection{Joint Floquet on $(t, b)$}

Module \texttt{floquet\_mobius} computes Floquet quasi-energies on
the joint time-circle $\times$ $\mathbb{Z}_3$ branch space. The
three-level Hamiltonian
\begin{equation}
H_{\text{static}} = \text{diag}(e_0, e_1, e_2)
  + g \cdot (\,|0\rangle\langle 1| + |1\rangle\langle 2| + |2\rangle\langle 0| + h.c.\,)
\end{equation}
is driven by $V(t) = A \sin(\Omega t) (S + S^\dagger)$, where $S$
is the $\mathbb{Z}_3$ cycle-shift. Quasi-energies satisfy
$\mathbb{Z}_3$ cyclic-permutation symmetry; resonances open at
$\Omega = e_b - e_{b'}$.

The \texttt{floquet\_engineering} extension module maps the
(amplitude, frequency) plane and identifies the CTC stability gap

\begin{equation}
\Delta_{\text{stable}}(\text{amp},\,\Omega)
  = \min_{i \ne j} |\epsilon_i - \epsilon_j| / \Omega,
\end{equation}
quantifying drive-induced stabilisation or destabilisation of CTC
bands.

\subsection{Tipler-sinusoid discrete-scale invariance and cascade-DSI}

The Tipler sinusoid $F(r) = \cos(\alpha \ln r/R + \delta)$ has
zero set $\{r_n = r_0 e^{\pi n / \alpha}\}$: a geometric progression
with ratio
\begin{equation}
\rho = e^{\pi / \alpha},
\end{equation}
verified to machine precision in \texttt{base\_tipler\_sinusoid\_is\_dsi}.
This is the textbook DSI signature of a log-periodic precursor
(Sornette 1998). A multi-cylinder cascade
\begin{equation}
F_{\text{cascade}}(r) = \sum_{k=0}^{L-1} A_k \cos(\alpha_k \ln r/R + \delta_k),
\quad \alpha_k = \alpha_0 \lambda^k,
\end{equation}
extends to non-trivial box-counting dimension (measured $> 0.3$ in
\texttt{cascade\_box\_dimension}). \texttt{dsi\_observables} ports
this machinery to data analysis: log-periodic precursor fits and
Lomb--Scargle log-frequency search.

\section{Back-reaction and self-consistency}

The \texttt{back\_reaction} module wires \texttt{hadamard\_offtrace}
and \texttt{newton\_kantorovich} into a composite back-reaction
residual
\begin{equation}
\mathcal{R}(\delta) = w_T \sum_r \|\langle T_{\mu\nu} \rangle(s_1; r)\|
  + w_L \sum_r |L_{\text{pair}}(\delta; r)|.
\end{equation}
For the matched pair (identical $a, R, A$), the residual is
minimised at $\delta = \pi$, matching the v0.11 CTC log-measure
result. This is \emph{not} a full Einstein-coupled back-reaction
solve (the spatial metric is not updated between iterations); it is
a static residual minimisation. Full back-reaction requires an
ADM-coupled solver, which is enabled by the \texttt{adm\_export}
module providing initial-data hand-off for Einstein Toolkit / similar
NR codes.

\section{D-CTC and quantum information}

Module \texttt{d\_ctc} implements Deutsch's 1991 CTC fixed-point
equation
\begin{equation}
\rho_{\text{CTC}} = \text{Tr}_{\text{CR}}
  \big[ U (\sigma_{\text{CR}} \otimes \rho_{\text{CTC}}) U^\dagger \big]
\end{equation}
on the $\mathbb{Z}_3$ branch space as the natural 3-dimensional
CTC register. The cyclic-shift unitary
$U = R_{\text{CR}} \otimes S$ admits the maximally-mixed CTC state
as a Banach fixed point; the iteration converges to machine
precision in tens of iterations.

This puts \textsc{Systroph\=e} in conversation with the
Aaronson--Watrous (2008) PSPACE result on D-CTC complexity: the
underlying spacetime geometry is here concretely instantiated.

\section{Validation modules for previously speculative items}

Six items previously catalogued as ``open speculative'' in
\texttt{docs/INTERPRETATIONS.md} have been promoted to concrete
calculations in v0.15.0. Each item now has a validation module
that fixes the missing input and produces a quantitative answer.

\paragraph{I.1: $\mathbb{Z}_3$ cover as wormhole throat
(\texttt{wormhole\_throat.py}).} The LP angular metric admits an
effective Morris-Thorne shape function $b_{\text{eff}}(r) = r - L(r) / r$.
The throat condition $b_{\text{eff}}(r_t) = r_t$ is exactly $L(r_t) = 0$,
i.e. the CTC band boundary. Flaring-out $b'(r_t) < 1$ holds at the
outer boundary of each CTC band.
\textbf{Verdict: conditional yes;} specific gluing constructed.

\paragraph{I.2: Casimir vs.\ exotic matter
(\texttt{exotic\_matter\_accounting.py}).} The Morris-Thorne exotic-
matter requirement is $|\rho_{\text{exotic}}|(r_t) = b'(r_t) / (8\pi r_t^2)$.
Equating this to the Casimir budget $|\rho_C| = \pi^2 / (720 d^4)$
gives the required plate separation
$d_{\text{req}} = (\pi^3 r_t^2 / (90 b'(r_t)))^{1/4}$. For $r_t \sim 1$
in natural units, $d_{\text{req}} / r_t \sim 0.7$: the cavity does
not geometrically fit within the throat.
\textbf{Verdict: quantitative NO.} The Casimir-replaces-exotic-matter
claim is falsified by direct dimensional accounting.

\paragraph{I.3: Self-consistent $\delta$
(\texttt{chronology\_protection.py}).} Newton-Kantorovich on the
composite back-reaction residual
$F(\delta) = w_T \sum_r |\langle T_{\mu\nu} \rangle|
+ w_L \sum_r |L_{\text{pair}}|$
from generic initial deltas converges to a $\delta$ in the
half-circle containing $\pi$ for the matched-pair case.
\textbf{Verdict: consistent} with chronology-protection conjecture.

\paragraph{I.4: M\"obius monodromy as U(1) Wilson loop
(\texttt{wilson\_loop.py}).} The flat connection
$A_\phi = (\gamma_{\text{eff}} + 2\pi b/3) / (2\pi)$ has Wilson loop
$W = \exp(i (\gamma_{\text{eff}} + 2\pi b / 3))$, reproducing the
$\mathbb{Z}_3$ branch phase. Field strength $F = dA = 0$ everywhere;
Chern number $c_1 = 0$. The non-trivial content is the holonomy
classifying $\pi_1(S^1) = \mathbb{Z}$.
\textbf{Verdict: concrete connection constructed.}

\paragraph{I.5: DCE photon flux (\texttt{dynamical\_casimir.py}).}
Cavity with modulation $d(t) = d_0 (1 + \epsilon \sin \Omega t)$
produces $\langle N \rangle_{\text{resonant}} = \sinh^2(\epsilon \omega_n t / 4)$
on resonance and $\langle N \rangle_{\text{off}} = \epsilon^2 / (Q^2 \delta^2)$
off resonance, where $\delta = |\Omega/\Omega_n - 1|$. The regime
selector picks the right formula by detuning vs.\ linewidth $1/Q$.
\textbf{Verdict: testable for any} $(Q, \epsilon, \delta)$.

\paragraph{I.6: Natural QFT vacuum (\texttt{vacuum\_states.py}).}
Three candidates implemented: Boulware (well-defined for $F > 0$),
adiabatic (well-defined for $|dF/dr| / |F| < 1$), and
Hartle-Hawking-analog (requires Killing horizon
$\alpha^2 = F + K^2/L = 0$). For supercritical LP, the chronology
horizon ($F = 0$) is generically \emph{not} a Killing horizon since
$K^2 / L \neq 0$ at $F = 0$.
\textbf{Verdict: NO.} The natural HH-analog vacuum does \emph{not}
exist on the generic supercritical LP exterior; the cavity-Casimir
state cannot be the natural vacuum in the standard QFTCS sense.
Use Boulware or adiabatic instead.

Two of the six items (I.2 and I.6) are \emph{negative} results: some
popular informal claims about Systrophe-style geometries are falsified
by direct calculation. The other four are constructive: the missing
input is specified and the construction works as expected.

\section{Summary and outlook}

\textsc{Systroph\=e} v0.15 provides the smallest open-source toolkit
known to the author that (i) constructs an exact closed-form CTC
spacetime, (ii) computes renormalised quantum-field-theoretic
observables on it, and (iii) supports both an experimental analog
(BEC vortices) and a quantum-information substrate (D-CTC).

\paragraph{Tests.} The package now contains 406 passing tests
covering all numerical claims in this paper. The full source is
at \texttt{github.com/Zynerji/systrophe}; license is MIT.

\paragraph{Next steps.} (a) BEC-vortex experimental design
collaboration (\texttt{CONTACTS.md}, Tier 1); (b) full
back-reaction NR run via ADM hand-off; (c) extension to the
multi-cylinder cascade DSI as a model of analog-cosmological
log-periodic precursors.

\section*{Acknowledgements}

This project was developed as a sustained dialogue between author
and the Anthropic ``Claude'' coding assistant; the codebase reflects
their joint authorship via per-commit \texttt{Co-Authored-By}
attribution. All physical claims and the experimental design are
the author's.

\bibliographystyle{plain}

\begin{thebibliography}{99}
\bibitem{Knopp2026Systrophe}
Knopp, C., \textit{Systroph\=e: A co-rotating Tipler-cylinder pair as a
tunable time-travel harness},
\texttt{github.com/Zynerji/systrophe}, 2026.

\bibitem{Deutsch1991}
Deutsch, D., \textit{Quantum mechanics near closed timelike curves},
Phys. Rev. D \textbf{44} (1991) 3197.

\bibitem{AaronsonWatrous2009}
Aaronson, S., Watrous, J., \textit{Closed timelike curves make
quantum and classical computing equivalent},
Proc. Roy. Soc. A \textbf{465} (2009) 631.

\bibitem{Unruh1981}
Unruh, W. G., \textit{Experimental black-hole evaporation?},
Phys. Rev. Lett. \textbf{46} (1981) 1351.

\bibitem{BrownMaclay1969}
Brown, L. S., Maclay, G. J.,
\textit{Vacuum stress between conducting plates: An image solution},
Phys. Rev. \textbf{184} (1969) 1272.

\bibitem{CallanHarvey1985}
Callan, C. G., Harvey, J. A., \textit{Anomalies and fermion zero modes
on strings and domain walls},
Nucl. Phys. B \textbf{250} (1985) 427.

\bibitem{Sornette1998}
Sornette, D., \textit{Discrete-scale invariance and complex dimensions},
Phys. Rep. \textbf{297} (1998) 239.

\bibitem{Christensen1976}
Christensen, S. M., \textit{Vacuum expectation value of the stress
tensor in an arbitrary curved background: The covariant point-separation
method}, Phys. Rev. D \textbf{14} (1976) 2490.

\bibitem{Steinhauer2019}
Steinhauer, J., \textit{Observation of thermal Hawking radiation and
its temperature in an analogue black hole}, Nature \textbf{569} (2019)
688.

\bibitem{Visser1996}
Visser, M., \textit{Lorentzian Wormholes: From Einstein to Hawking},
AIP Press, 1996.

\end{thebibliography}

\end{document}
